\documentclass[]{beamer}

\usetheme{metropolis}
\metroset{progressbar=frametitle}
\usefonttheme[onlymath]{serif}

\usepackage[none]{hyphenat}
\usepackage{hyperref}
\usepackage{amsmath}
\usepackage{amsfonts}
\usepackage{amssymb}
\usepackage[normalem]{ulem}
\usepackage{graphicx}
\usepackage{tikz}
\usepackage{minted}
\subtitle{Anomaly Detection and Robust Methods}
\author{Maxim Borisyak}

\institute{Constructor University Bremen}

\usepackage{algorithm}
\usepackage{booktabs}
\usepackage{algpseudocode}
\usepackage{setspace}
\usepackage{framed}

\DeclareMathOperator*{\E}{\mathbb{E}}
\DeclareMathOperator*{\var}{\mathbb{D}}
\newcommand\D[1]{\var\left[ #1 \right]}

\DeclareMathOperator*{\argmin}{\mathrm{arg\,min}}
\DeclareMathOperator*{\argmax}{\mathrm{arg\,max}}




\begin{document}



\title{Machine Learning and Data Mining}

\begin{frame}[plain]
\titlepage
\end{frame}

\section{Anomaly Detection}

\begin{frame}[fragile]
\frametitle{What is an anomaly?}
\textbf{Anomaly} (outlier): an observation that deviates markedly from the majority of the data.

\textbf{Applications}: fraud detection, network intrusion, fault detection, medical diagnosis.

\textbf{Types}:
\begin{itemize}
  \item \textbf{point}: a single observation is anomalous (e.g.\ a \$10\,000 transaction among \$20--\$200 ones);
  \item \textbf{contextual}: anomalous in context but normal globally (e.g.\ $30^\circ$C in January);
  \item \textbf{collective}: a group of individually normal observations is anomalous together.
\end{itemize}

\end{frame}

\begin{frame}[fragile]
\frametitle{Problem formulations}
\textbf{Unsupervised}: only unlabeled data; score each point by deviation from the bulk.

\textbf{Semi-supervised (one-class)}: train on clean normal data; flag deviations at test time.

\textbf{Supervised}: labeled normal and anomalous examples; standard binary classification, but highly imbalanced.

\end{frame}

\begin{frame}[fragile]
\frametitle{Statistical approach: $z$-score}
\begin{equation*}
z_i = \frac{x_i - \mu}{\sigma}
\end{equation*}

Flag $x_i$ as outlier if $|z_i| > \tau$ (commonly $\tau = 3$).

\textbf{Weakness}: $\mu$ and $\sigma$ are themselves distorted by outliers --- extreme values inflate $\sigma$, masking the anomalies we seek.

\end{frame}

\begin{frame}[fragile]
\frametitle{Statistical approach: IQR rule}
Let $\mathrm{IQR} = Q_3 - Q_1$. Flag $x_i$ if:
\begin{equation*}
x_i < Q_1 - 1.5\,\mathrm{IQR} \quad \text{or} \quad x_i > Q_3 + 1.5\,\mathrm{IQR}
\end{equation*}

\begin{itemize}
  \item quartiles are unaffected by extreme values;
  \item threshold $1.5$ corresponds to $\approx \pm 2.7\sigma$ under Gaussianity.
\end{itemize}

\end{frame}

\begin{frame}[fragile]
\frametitle{Distance-based outlier detection}
\textbf{$k$-NN score}: $\mathrm{score}(x) = d_k(x)$, the distance to the $k$-th nearest neighbor.

\begin{itemize}
  \item no distributional assumptions;
  \item $O(n^2)$ pairwise distances --- slow for large $n$.
\end{itemize}

\textbf{Local Outlier Factor (LOF)}: compares local density of $x$ to that of its neighbors.

\begin{equation*}
\mathrm{LOF}_k(x) = \frac{1}{k}\sum_{o \in \mathcal{N}_k(x)} \frac{\mathrm{lrd}_k(o)}{\mathrm{lrd}_k(x)}
\end{equation*}

$\mathrm{LOF} \approx 1$: normal; $\mathrm{LOF} \gg 1$: anomalous.

\end{frame}

\section{One-Class SVM}

\begin{frame}[fragile]
\frametitle{One-class SVM}
\textbf{Setting}: only normal training data $\{x_1,\dots,x_n\}$.

\textbf{Goal}: find a hypersphere in feature space enclosing most of the data; points outside are anomalies.

\begin{equation*}
\min_{w,\, \xi,\, \rho} \;\frac{1}{2}\|w\|^2 - \rho + \frac{1}{\nu n}\sum_{i=1}^{n}\xi_i
\quad \text{s.t.} \quad
\langle w,\, \phi(x_i)\rangle \geq \rho - \xi_i,\;\; \xi_i \geq 0
\end{equation*}

\begin{itemize}
  \item $\nu \in (0,1]$: upper bound on the outlier fraction and lower bound on support vectors;
  \item $\rho$: offset from the origin (margin).
\end{itemize}

\end{frame}

\begin{frame}[fragile]
\frametitle{One-class SVM: decision function}
\begin{equation*}
f(x) = \mathrm{sign}\!\left(\sum_{i} \alpha_i\, K(x_i, x) - \rho\right)
\end{equation*}

Common choice: RBF kernel $K(x, x') = \exp(-\gamma \|x - x'\|^2)$.

\textbf{Pros}: works in high dimensions; $\nu$ gives direct control over outlier fraction.

\textbf{Cons}: $O(n^2)$--$O(n^3)$ training; sensitive to kernel choice; training data must be clean.

\end{frame}

\section{Isolation Forest}

\begin{frame}[fragile]
\frametitle{Isolation Forest: key idea}
\textbf{Principle} (Liu, Ting, Zhou 2008): anomalies are few and different --- they are easier to isolate.

Build random binary trees by repeatedly splitting on a random feature at a random threshold. A point is isolated when it reaches a leaf alone.

\begin{itemize}
  \item normal point: buried in a dense region, needs many splits;
  \item anomalous point: in a sparse region, isolated in few splits.
\end{itemize}

\end{frame}

\begin{frame}[fragile]
\frametitle{Isolation Forest: score}
Build $t$ trees on subsamples of size $\psi$. \textbf{Anomaly score} (normalized to $[0,1]$):

\begin{equation*}
s(x,\psi) = 2^{-\,\bar{h}(x) / c(\psi)}, \qquad c(\psi) = 2H(\psi - 1) - \frac{2(\psi-1)}{\psi}
\end{equation*}

where $\bar{h}(x)$ is the mean path length over the forest and $c(\psi)$ is the expected BST depth.

$s \approx 1$: anomaly; $s \approx 0.5$: normal.

\textbf{Complexity}: $O(t\,\psi\log\psi)$ --- linear in $n$, no distance computation.

\end{frame}

\begin{frame}[fragile]
\frametitle{Comparing anomaly detectors}
\begin{center}
{\footnotesize\setlength\tabcolsep{4pt}
\begin{tabular}{lllll}
\toprule
\textbf{Method} & \textbf{Complexity} & \textbf{High-dim} & \textbf{Density} & \textbf{Interpretable} \\
\midrule
$z$-score / IQR  & $O(n)$             & No  & No  & Yes \\
$k$-NN / LOF     & $O(n^2)$           & No  & Yes & Partial \\
One-class SVM    & $O(n^2)$--$O(n^3)$ & Yes & Yes & No \\
Isolation Forest & $O(n)$             & Yes & No  & No \\
\bottomrule
\end{tabular}}
\end{center}

\end{frame}

\section{Robust Regression}

\begin{frame}[fragile]
\frametitle{Why OLS fails}
\textbf{OLS} minimizes $\sum_i r_i^2$ where $r_i = y_i - x_i^\top\beta$.

The squared loss penalizes large residuals quadratically --- a single outlier can dominate the fit.

\textbf{Breakdown point}: the fraction of outliers an estimator can tolerate.

\begin{itemize}
  \item OLS: $1/n \to 0$ (a single outlier suffices to corrupt the estimate).
\end{itemize}

\end{frame}

\begin{frame}[fragile]
\frametitle{Robust loss functions}
Replace $r^2$ with $\rho(r)$ that grows slowly for large $|r|$:

\begin{center}
\begin{tabular}{lll}
\toprule
\textbf{Loss} & $\rho(r)$ & $\psi(r) = \rho'(r)$ \\
\midrule
OLS       & $r^2$                           & $2r$ (unbounded) \\
Huber     & $r^2/2$ or $\delta|r|-\delta^2/2$ & bounded \\
Cauchy    & $\log(1 + r^2/\sigma^2)$        & bounded \\
Tukey     & piecewise cubic                 & redescending (zero for large $|r|$) \\
\bottomrule
\end{tabular}
\end{center}

\end{frame}

\begin{frame}[fragile]
\frametitle{Huber loss and IRLS}
\begin{equation*}
\rho_\delta(r) =
\begin{cases}
r^2/2 & |r| \leq \delta \\
\delta\,|r| - \delta^2/2 & |r| > \delta
\end{cases}
\end{equation*}

Quadratic near zero, linear in the tails. $\delta$ is typically set via MAD of residuals.

\textbf{IRLS}: minimizing $\sum_i\rho(r_i)$ is equivalent to weighted least squares with adaptive weights $w_i = \psi(r_i)/r_i$:

\begin{equation*}
\beta^{(t+1)} = \left(X^\top W^{(t)} X\right)^{-1} X^\top W^{(t)} y
\end{equation*}

\end{frame}

\begin{frame}[fragile]
\frametitle{RANSAC}
\textbf{Random Sample Consensus}: model-agnostic robust fitting.

\begin{framed}
\begin{spacing}{1.4}
\begin{algorithmic}[1]
  \Repeat
    \State sample $s$ points; fit model $\hat{\beta}$
    \State find inliers $\mathcal{I} = \{i : |r_i(\hat{\beta})| < \tau\}$
    \State keep $\hat{\beta}$ if $|\mathcal{I}|$ is largest so far
  \Until{$T$ iterations}
  \State refit on all inliers $\mathcal{I}^*$
\end{algorithmic}
\end{spacing}
\end{framed}

Required iterations to succeed with probability $1-\eta$:
$T \geq \log\eta \,/\, \log[1-(1-\epsilon)^s]$, where $\epsilon$ = outlier fraction.

\end{frame}

\begin{frame}[fragile]
\frametitle{Theil--Sen estimator}
Slope estimated as the \textbf{median pairwise slope}:

\begin{equation*}
\hat{\beta}_1 = \mathrm{median}_{i < j}\; \frac{y_j - y_i}{x_j - x_i}, \qquad
\hat{\beta}_0 = \mathrm{median}_i(y_i - \hat{\beta}_1 x_i)
\end{equation*}

\begin{itemize}
  \item breakdown point $\approx 29\%$;
  \item $O(n^2)$ naively, $O(n\log n)$ with efficient median algorithms.
\end{itemize}

\end{frame}

\section{Handling Missing Values}

\begin{frame}[fragile]
\frametitle{Types of missingness}
\textbf{MCAR} (Missing Completely At Random): $P(\text{miss}) $ independent of all values.
\begin{itemize}
  \item complete-case analysis is unbiased.
\end{itemize}

\textbf{MAR} (Missing At Random): $P(\text{miss} \mid x_\text{obs}, x_\text{miss}) = P(\text{miss} \mid x_\text{obs})$.
\begin{itemize}
  \item bias unless conditioning on observed variables.
\end{itemize}

\textbf{MNAR} (Missing Not At Random): $P(\text{miss})$ depends on the missing value itself.
\begin{itemize}
  \item requires modeling the missingness mechanism explicitly.
\end{itemize}

\end{frame}

\begin{frame}[fragile]
\frametitle{Why missingness type matters}
\begin{center}
\begin{tabular}{lp{3.2cm}p{3.2cm}}
\toprule
\textbf{Mechanism} & \textbf{Complete-case} & \textbf{Simple imputation} \\
\midrule
MCAR & Unbiased           & Unbiased \\
MAR  & Biased             & Unbiased if conditioned correctly \\
MNAR & Biased             & Biased \\
\bottomrule
\end{tabular}
\end{center}

\vspace{3mm}
\textbf{Complete-case analysis}: drop rows with any missing value --- simple but wasteful and biased unless MCAR.

\end{frame}

\begin{frame}[fragile]
\frametitle{Imputation strategies}
\textbf{Univariate} (per-column): mean / median / mode / constant.
\begin{itemize}
  \item distorts variance; ignores correlations.
\end{itemize}

\textbf{$k$-NN imputation}: fill from weighted mean of $k$ nearest complete-case neighbors.

\textbf{Regression imputation}: regress $x_j$ on all other features; predict missing entries.

\textbf{MICE} (Multiple Imputation by Chained Equations): iterate regression imputation across features.

\end{frame}

\begin{frame}[fragile]
\frametitle{Matrix completion}
\textbf{Low-rank assumption}: $X \approx Z$ with $\mathrm{rank}(Z) \ll d$.

\begin{equation*}
\min_{Z}\;\sum_{(i,j)\in\Omega} \left(X_{ij} - Z_{ij}\right)^2 + \lambda\|Z\|_*
\end{equation*}

Nuclear norm $\|Z\|_* = \sum_j \sigma_j(Z)$ is the convex surrogate for rank. Same machinery as collaborative filtering.

\textbf{Native support}: XGBoost, LightGBM, CatBoost send missing-value samples to the child that minimizes loss --- no imputation needed.

\end{frame}

\begin{frame}[fragile]
\frametitle{Summary: anomaly detection}
\begin{itemize}
  \item $z$-score / IQR: simple, univariate;
  \item LOF: density-aware, handles varying density;
  \item one-class SVM: kernel boundary, high-dimensional;
  \item Isolation Forest: fast ($O(n)$), scalable.
\end{itemize}

\end{frame}

\begin{frame}[fragile]
\frametitle{Summary: robust regression}
\begin{itemize}
  \item Huber + IRLS: smooth, adaptive down-weighting;
  \item RANSAC: model-agnostic, handles high outlier fractions;
  \item Theil--Sen: 29\% breakdown point, simple regression.
\end{itemize}

\end{frame}

\begin{frame}[fragile]
\frametitle{Summary: missing values}
\begin{itemize}
  \item identify mechanism: MCAR / MAR / MNAR;
  \item MCAR: simple imputation; MAR: MICE or matrix completion;
  \item MNAR: model the missingness itself;
  \item tree models handle missing values natively.
\end{itemize}

\end{frame}




\end{document}
